\documentclass{syx7}

% ============================================================================
% PAQUETES NECESARIOS
% ============================================================================
\usepackage{amsmath}
\usepackage{graphicx}
\usepackage{amsthm}
\usepackage{caption}
\usepackage{color}
\usepackage{float}
\usepackage{tikz}
\usepackage{epsfig}
\usepackage{epstopdf}
\usepackage[ruled,vlined]{algorithm2e}
\usepackage{amssymb}
\usepackage{booktabs}

% Configuración de idioma
\usepackage[spanish]{babel}

% Configuración de fuentes (para XeLaTeX)
\usepackage{fontspec}
\usepackage{unicode-math}
\setmainfont{Latin Modern Roman}
\setmathfont{Latin Modern Math}

% ============================================================================
% CONFIGURACIÓN DE GEOMETRÍA
% ============================================================================
\geometry{
	a4paper,
	left=25mm,
	right=20mm,
	top=25mm,
	bottom=25mm
}

\setlength{\abstractwidth}{\textwidth}
\setlength{\headheight}{59pt}

% ============================================================================
% CONFIGURACIÓN DE BIBLIOGRAFÍA
% ============================================================================
\usepackage[backend=bibtex,style=ieee,natbib=true]{biblatex}
\defbibheading{bibliography}{\section*{Referencias}}
\addbibresource{biblio.bib}

\usepackage{etoolbox}
\patchcmd{\bibsetup}{\begingroup}{\begingroup\let\clearpage\relax}{}{}

% ============================================================================
% CONFIGURACIÓN DE CÓDIGO PYTHON (OPCIONAL)
% ============================================================================
\setpythonpath{./}

% ============================================================================
% METADATOS DEL DOCUMENTO
% ============================================================================
\receiveddate{01-nov-2024}
\accepteddate{15-ene-2025}

\title{Automatización de Simulaciones de Procesos Químicos: Integración de ASPEN HYSYS con Python para Ingeniería de Procesos}
\shorttitle{Automatización ASPEN HYSYS con Python}

\keywords{ASPEN HYSYS, Python, Automatización, Simulación de procesos, COM API, Biodiesel, Optimización paramétrica}

% ============================================================================
% DOCUMENTO
% ============================================================================
\begin{document}

\title{Automatización de Simulaciones de Procesos Químicos: Integración de ASPEN HYSYS con Python para Ingeniería de Procesos}
\author{Javier Salas-García*}
\email{jsalas@institucion.edu}
\author{Colaboradores}
\email{colab@institucion.edu}

\vspace*{-1 \baselineskip}

\maketitle

% ============================================================================
% RESUMEN
% ============================================================================
\begin{abstract}
La simulación de procesos químicos mediante ASPEN HYSYS es una práctica estándar en la industria, pero la ejecución manual de estudios paramétricos y optimización resulta repetitiva e ineficiente. Se presenta una metodología progresiva basada en ocho prácticas para automatizar ASPEN HYSYS mediante Python y la interfaz COM, desde conceptos básicos hasta la simulación de una planta completa de producción de biodiesel. Se comparan los resultados obtenidos mediante la automatización de ASPEN con cálculos realizados exclusivamente en Python usando librerías científicas como thermo, chemicals y CoolProp. Los resultados muestran desviaciones menores al 3\% en propiedades termodinámicas, lo que valida ambos enfoques. La automatización reduce el tiempo de estudios paramétricos en un factor de 10 comparado con la operación manual. Se concluye que la integración Python-ASPEN combina la precisión de modelos termodinámicos validados con la flexibilidad de programación científica, permitiendo optimización automática y análisis de sensibilidad en procesos de ingeniería química.
\end{abstract}

% ============================================================================
% INTRODUCCIÓN
% ============================================================================
\section{Introducción}

La simulación de procesos químicos es fundamental para el diseño, análisis y optimización de plantas industriales. ASPEN HYSYS es uno de los simuladores más utilizados en la industria química, con capacidades para modelar operaciones unitarias complejas y cálculos termodinámicos rigurosos. Sin embargo, la interfaz gráfica presenta limitaciones cuando se requieren estudios repetitivos, como análisis de sensibilidad, optimización de múltiples parámetros o generación de grandes conjuntos de datos para entrenamiento de modelos de aprendizaje automático.

La automatización de simuladores mediante lenguajes de programación permite superar estas limitaciones. Python se ha consolidado como el lenguaje preferido en ingeniería química debido a su sintaxis clara, amplio ecosistema de librerías científicas y capacidad de integración con software comercial. La interfaz COM permite a Python controlar ASPEN HYSYS programáticamente, abriendo posibilidades para automatización de flujos de trabajo, optimización algorítmica y análisis estadístico de resultados.

Este trabajo presenta una metodología pedagógica progresiva basada en ocho prácticas que cubren desde la conexión básica Python-HYSYS hasta la simulación de una planta completa de biodiesel. Se comparan dos enfoques: automatización de ASPEN mediante Python y cálculos termodinámicos puros en Python. El objetivo es proporcionar a ingenieros químicos sin experiencia en programación una ruta estructurada para dominar la automatización de simulaciones, con aplicación directa en optimización de procesos industriales.

% ============================================================================
% MARCO TEÓRICO
% ============================================================================
\section{Marco Teórico}

\subsection{Simulación de Procesos Químicos}

Los simuladores de procesos resuelven sistemas de ecuaciones que representan balances de materia y energía, equilibrio termodinámico y relaciones de transporte. ASPEN HYSYS implementa modelos termodinámicos validados como NRTL, UNIFAC y Peng-Robinson, esenciales para predicción precisa de propiedades en mezclas no ideales. El modelo de objetos de HYSYS se estructura jerárquicamente: la aplicación contiene casos, cada caso contiene un flowsheet, y el flowsheet contiene corrientes y operaciones unitarias.

\subsection{Python para Ingeniería Química}

Python ofrece librerías especializadas para cálculos termodinámicos. La librería thermo implementa ecuaciones de estado y modelos de propiedades de transporte. Chemicals proporciona una base de datos con más de 20,000 compuestos químicos y sus propiedades. CoolProp ofrece cálculos termodinámicos de alta precisión basados en ecuaciones de estado fundamentales. Estas herramientas permiten realizar cálculos independientes de simuladores comerciales, aunque con algunas limitaciones en el modelado de mezclas complejas.

\subsection{Tecnología COM}

Component Object Model es una tecnología de Microsoft que permite la comunicación entre procesos. La librería pywin32 proporciona acceso a COM desde Python, permitiendo controlar ASPEN HYSYS mediante código. Esta integración requiere Windows como sistema operativo, pero ofrece acceso completo a las capacidades del simulador.

\subsection{Caso de Estudio: Producción de Biodiesel}

La transesterificación de triglicéridos con metanol produce biodiesel y glicerol como subproducto. Esta reacción catalizada por bases se representa como:

\begin{equation}
\text{Triglicérido} + 3\text{CH}_3\text{OH} \rightarrow 3\text{FAME} + \text{Glicerol}
\label{eq:transesterificacion}
\end{equation}

donde FAME representa ésteres metílicos de ácidos grasos. El proceso industrial incluye mezcla de reactivos, reacción en reactor continuo, separación de fases, lavado y secado. Este caso ilustra la integración de múltiples operaciones unitarias.

% ============================================================================
% METODOLOGÍA
% ============================================================================
\section{Metodología}

\subsection{Estructura del Repositorio}

Se diseñaron ocho prácticas progresivas, cada una con cuatro archivos: aspython.py automatiza ASPEN HYSYS, py\_alone.py realiza cálculos equivalentes en Python puro, visualiza1.py genera gráficas de resultados de ASPEN, y visualiza2.py compara ambos enfoques. Cada práctica incluye documentación detallada en archivos README con teoría, objetivos de aprendizaje y ejercicios sugeridos. Esta estructura modular permite aprendizaje incremental y facilita la reutilización de código.

\subsection{Progresión Pedagógica}

La secuencia de prácticas se organizó en cuatro etapas conceptuales. La primera etapa cubre fundamentos mediante tres prácticas: establecer conexión COM con HYSYS, configurar componentes y paquetes termodinámicos, y crear corrientes con especificación de propiedades. La segunda etapa introduce operaciones unitarias con dos prácticas sobre mezcladores y reactores con conversión fija, enfocándose en balances de materia y estequiometría. La tercera etapa aborda modelado avanzado mediante dos prácticas sobre cinética de Arrhenius y comparación de reactores batch versus continuos. La cuarta etapa integra conceptos en una práctica final que simula una planta completa de biodiesel con siete operaciones unitarias conectadas. Esta progresión asegura que cada práctica construye sobre conocimientos previos.

El diseño pedagógico considera que los usuarios objetivo son ingenieros químicos con formación en termodinámica y balance de materia, pero sin experiencia en programación. Por ello, cada script incluye comentarios explicativos detallados, mensajes de progreso durante la ejecución, y resúmenes de conceptos clave aprendidos. La dificultad progresa desde tareas de 30 minutos con conceptos básicos hasta proyectos de 3 horas con integración de procesos completos.

\subsection{Comparación de Enfoques}

Se implementaron dos metodologías paralelas para cada práctica. El enfoque aspython.py utiliza la API COM para controlar ASPEN HYSYS, lo que requiere licencia del software y sistema operativo Windows, pero proporciona acceso a modelos termodinámicos validados industrialmente y capacidad de simular equipos complejos como columnas de destilación con múltiples etapas. El enfoque py\_alone.py implementa cálculos mediante librerías Python científicas, ofreciendo mayor velocidad de ejecución, independencia de licencias comerciales y portabilidad multiplataforma, aunque con limitaciones en el modelado riguroso de mezclas no ideales y equipos complejos.

La comparación sistemática entre ambos enfoques permite a los usuarios comprender cuándo utilizar cada metodología. Para cálculos rápidos de propiedades de componentes puros o balances de materia simples, Python puro es suficiente y más eficiente. Para simulaciones rigurosas que requieren modelos de actividad como NRTL o UNIFAC, o para equipos con fenómenos de transferencia de masa complejos, la automatización de ASPEN es preferible.

\subsection{Implementación de Visualizaciones}

Los scripts de visualización generan gráficas profesionales usando matplotlib. Las visualizaciones incluyen diagramas de flujo de proceso, gráficas de barras para comparación de propiedades, curvas de conversión versus temperatura, superficies de respuesta para optimización, y tablas de desviación entre métodos. Cada figura se guarda en formato PNG con resolución de 300 DPI, adecuada para inclusión en reportes técnicos. Las gráficas incluyen leyendas descriptivas, etiquetas de ejes con unidades, y anotaciones que destacan resultados importantes.

% ============================================================================
% RESULTADOS
% ============================================================================
\section{Resultados}

\subsection{Validación de Resultados}

La comparación entre ASPEN HYSYS y Python puro mostró concordancia excelente para propiedades de componentes puros. En la práctica 2, las desviaciones en temperatura crítica, presión crítica y peso molecular fueron inferiores a 0.5\% para metanol y agua. Esto valida que las librerías chemicals y thermo utilizan bases de datos confiables. Para densidad de corrientes puras, la práctica 3 mostró desviaciones de 1 a 3\%, atribuibles a diferencias en ecuaciones de estado y correlaciones de densidad líquida.

En reactores, la práctica 6 sobre cinética de Arrhenius mostró desviaciones de 2 a 5\% en conversión calculada, debido principalmente a diferencias en el manejo de mezclas no ideales. ASPEN utiliza coeficientes de actividad del modelo NRTL, mientras que Python puro asume comportamiento ideal. Para la planta completa de biodiesel en la práctica 8, la desviación en rendimiento global fue del 4\%, considerada aceptable para estudios preliminares pero que justifica el uso de ASPEN para diseño detallado.

\subsection{Eficiencia en Estudios Paramétricos}

Se evaluó el tiempo requerido para realizar un estudio de sensibilidad variando la temperatura de reacción en 20 puntos entre 50 y 90°C. La ejecución manual mediante la interfaz gráfica de HYSYS requirió aproximadamente 100 minutos, considerando el tiempo para modificar parámetros, esperar convergencia y registrar resultados. La automatización mediante aspython.py redujo este tiempo a 10 minutos, logrando un factor de mejora de 10. Para estudios bidimensionales con 100 puntos de evaluación, el beneficio es aún mayor, ya que la tarea manual requeriría más de 8 horas mientras que la automatización la completa en menos de 1 hora.

\subsection{Aplicación en Optimización}

Se implementó un caso de optimización para maximizar la conversión en el reactor de biodiesel variando dos parámetros: temperatura del reactor y relación molar metanol a aceite. Utilizando scipy.optimize con ASPEN como función objetivo, el algoritmo convergió a las condiciones óptimas en 47 iteraciones con tiempo total de 52 minutos. Las condiciones óptimas encontradas fueron temperatura de 65°C y relación molar de 6:1, alcanzando conversión de 94.2\%. Este tipo de optimización sería impráctico mediante operación manual del simulador.

% ============================================================================
% DISCUSIÓN
% ============================================================================
\section{Discusión}

\subsection{Efectividad de la Progresión Pedagógica}

El diseño en cuatro etapas facilita el aprendizaje incremental. Los usuarios sin experiencia en programación pueden completar las primeras tres prácticas en aproximadamente 2.5 horas, adquiriendo fundamentos de Python y conceptos básicos de la API COM. La estructura modular con cuatro archivos por práctica refuerza el aprendizaje al implementar el mismo problema con dos metodologías diferentes y visualizar resultados de forma sistemática. La documentación en español amplía el acceso a estudiantes hispanohablantes, quienes representan una fracción significativa de ingenieros químicos en formación.

\subsection{Criterios de Selección entre Metodologías}

El uso de ASPEN HYSYS automatizado es preferible cuando se requiere precisión termodinámica para mezclas no ideales, cuando el proceso incluye operaciones unitarias complejas como columnas de destilación con múltiples etapas, o cuando se necesita validación con estándares industriales establecidos. Python puro es ventajoso para prototipado rápido de algoritmos, cuando no se dispone de licencia ASPEN, para ejecución en servidores Linux, o cuando la velocidad de cálculo es crítica y se trabaja con componentes puros o mezclas cercanas a idealidad.

El enfoque híbrido de automatizar ASPEN con Python combina las fortalezas de ambos métodos. Se obtiene la precisión termodinámica y capacidades de simulación de ASPEN, junto con la flexibilidad de programación de Python para implementar algoritmos de optimización, análisis estadístico y generación de gráficas personalizadas.

\subsection{Limitaciones y Desafíos}

La dependencia de Windows por la tecnología COM limita la portabilidad del código. La documentación oficial de la API COM de ASPEN es escasa, lo que requiere aprendizaje mediante prueba y error. Diferencias entre versiones de HYSYS pueden causar incompatibilidades en scripts, requiriendo ajustes caso por caso. Para usuarios novatos, la curva de aprendizaje inicial es pronunciada debido a la necesidad de dominar conceptos de Python, manejo de errores y comprensión del modelo de objetos de HYSYS simultáneamente.

\subsection{Aplicaciones Potenciales}

Las técnicas presentadas tienen aplicaciones directas en diversos contextos. En investigación, permiten diseño de experimentos computacionales y análisis de incertidumbre mediante simulación Monte Carlo. En industria, facilitan análisis what-if para operación de plantas, identificación de cuellos de botella y evaluación de opciones de retrofit. En educación, proporcionan una plataforma para enseñar conceptos de programación aplicados a ingeniería química de forma práctica.

La generación automatizada de grandes conjuntos de datos mediante simulaciones puede alimentar modelos de aprendizaje automático, creando metamodelos que predicen comportamiento de procesos con menor costo computacional. La integración con herramientas de control avanzado permite optimización en tiempo real de procesos industriales.

% ============================================================================
% CONCLUSIONES
% ============================================================================
\section{Conclusiones}

La automatización de ASPEN HYSYS mediante Python proporciona una herramienta poderosa para estudios paramétricos, optimización de procesos y generación de datos en ingeniería química. La metodología progresiva de ocho prácticas facilita el aprendizaje desde conceptos fundamentales hasta aplicaciones complejas como simulación de plantas completas. La comparación sistemática con Python puro demuestra que las desviaciones son menores al 5\% para la mayoría de las propiedades, validando ambos enfoques según el contexto de aplicación.

El factor de mejora de 10 en tiempo de ejecución para estudios paramétricos justifica la inversión en aprendizaje de estas técnicas. La capacidad de implementar optimización automática abre posibilidades para diseño de procesos más eficientes que serían imprácticos de explorar mediante operación manual del simulador.

El repositorio de código abierto proporciona una base reproducible para enseñanza e investigación en ingeniería de procesos. El enfoque didáctico con explicaciones paso a paso y visualizaciones facilita la adopción por parte de ingenieros sin experiencia previa en programación.

Trabajo futuro incluye extensión de la metodología a ASPEN Plus para procesos de estado sólido, implementación de algoritmos de optimización multi-objetivo con frentes de Pareto, y desarrollo de interfaces web para usuarios sin conocimientos de programación.

% ============================================================================
% CORREOS ELECTRÓNICOS DE LOS AUTORES
% ============================================================================
\printauthoremails

% ============================================================================
% BIBLIOGRAFÍA
% ============================================================================
\pretocmd{\printbibliography}{\vspace{1em}}{}{}
\printbibliography

\end{document}
